\documentclass[11pt,a4paper]{article}
\usepackage[utf8]{inputenc}
\usepackage[spanish]{babel}
\usepackage{graphicx}
\usepackage{geometry}
\usepackage{caption}
\usepackage{subcaption}
\usepackage{hyperref}
\usepackage{float}
\geometry{top=2cm, bottom=2cm, left=2.5cm, right=2.5cm}

\title{Reporte de Experimento EMG}
\author{Ñandú - Sistema de Adquisición EMG}
\date{2026-09-22}

\begin{document}

\maketitle

\section{Configuración y Hardware}
\textbf{Baterías:} Alimentación por baterías de 9V (bajo ruido) \\
\textbf{Tierra:} Referencia GND en apófisis mastoides

\section{Armado y Electrodos}
Registro de superficie sEMG submáximal (Sujeto1)

\section{Ubicación de los Músculos}
\begin{itemize}
    \item \textbf{Canal 0:} masetero
    \item \textbf{Canal 1:} -
    \item \textbf{Canal 2:} orbi
\end{itemize}
Canal 0: masetero, Canal 1: -, Canal 2: orbi

\section{Protocolo y Secuencia}
Secuencia periódica de fonación vocálica guiada por metrónomo a 30 BPM

\section{Notas y Observaciones}
Reporte técnico integral compilado de forma automatizada por el sistema.

\newpage
\section{Patrones Musculares Promedio}
Comparativa de los patrones musculares promedio suavizados para las 5 vocales fonatorias:

\begin{figure}[H]
\centering
\begin{subfigure}{0.31\textwidth}
    \includegraphics[width=\textwidth]{/home/santiago/repositorios/Nandu_SistemadeAdqusicionEMG/EMG_desarrollo/base_de_datos_electrodos/2026-09-22/O_Prueba6_Sujeto1/patron_muscular_grabacion.png}
    \caption{Vocal O}
\end{subfigure}
\hfill
\begin{subfigure}{0.31\textwidth}
    \includegraphics[width=\textwidth]{/home/santiago/repositorios/Nandu_SistemadeAdqusicionEMG/EMG_desarrollo/base_de_datos_electrodos/2026-09-22/U_Prueba6_Sujeto1/patron_muscular_grabacion.png}
    \caption{Vocal U}
\end{subfigure}
\caption{Comparativa de patrones musculares promedio suavizados.}
\end{figure}

\newpage
\section{Análisis de Señales por Vocal}
A continuación se presentan los registros de 3 músculos del paper y las señales calibradas agrupadas vocal por vocal:

\subsection{Vocal O}
Se presentan a continuación las tomas registradas para la vocal \textbf{O}, ordenadas cronológicamente por número de prueba:

\subsubsection{Medición: O\_Prueba6\_Sujeto1}
\begin{figure}[H]
\centering
\includegraphics[width=0.88\textwidth]{/home/santiago/repositorios/Nandu_SistemadeAdqusicionEMG/EMG_desarrollo/base_de_datos_electrodos/2026-09-22/O_Prueba6_Sujeto1/plot_paper_combined.png}
\caption{Registro de 3 músculos de la vocal O - O\_Prueba6\_Sujeto1: línea temporal de ventanas y concatenación simétrica con ruido restado.}
\end{figure}

\begin{figure}[H]
\centering
\includegraphics[width=0.88\textwidth]{/home/santiago/repositorios/Nandu_SistemadeAdqusicionEMG/EMG_desarrollo/base_de_datos_electrodos/2026-09-22/O_Prueba6_Sujeto1/plot_calibrado_2026-09-22_O_Prueba6_Sujeto1.png}
\caption{Señales calibradas de la vocal O - O\_Prueba6\_Sujeto1: filtro Notch en 50 Hz, pasabanda y envolvente RMS.}
\end{figure}

\subsubsection{Medición: O\_Prueba7\_Sujeto1}
\begin{figure}[H]
\centering
\includegraphics[width=0.88\textwidth]{/home/santiago/repositorios/Nandu_SistemadeAdqusicionEMG/EMG_desarrollo/base_de_datos_electrodos/2026-09-22/O_Prueba7_Sujeto1/plot_paper_combined.png}
\caption{Registro de 3 músculos de la vocal O - O\_Prueba7\_Sujeto1: línea temporal de ventanas y concatenación simétrica con ruido restado.}
\end{figure}

\begin{figure}[H]
\centering
\includegraphics[width=0.88\textwidth]{/home/santiago/repositorios/Nandu_SistemadeAdqusicionEMG/EMG_desarrollo/base_de_datos_electrodos/2026-09-22/O_Prueba7_Sujeto1/plot_calibrado_2026-09-22_O_Prueba7_Sujeto1.png}
\caption{Señales calibradas de la vocal O - O\_Prueba7\_Sujeto1: filtro Notch en 50 Hz, pasabanda y envolvente RMS.}
\end{figure}

\subsubsection{Medición: O\_Prueba8\_Sujeto1}
\begin{figure}[H]
\centering
\includegraphics[width=0.88\textwidth]{/home/santiago/repositorios/Nandu_SistemadeAdqusicionEMG/EMG_desarrollo/base_de_datos_electrodos/2026-09-22/O_Prueba8_Sujeto1/plot_paper_combined.png}
\caption{Registro de 3 músculos de la vocal O - O\_Prueba8\_Sujeto1: línea temporal de ventanas y concatenación simétrica con ruido restado.}
\end{figure}

\begin{figure}[H]
\centering
\includegraphics[width=0.88\textwidth]{/home/santiago/repositorios/Nandu_SistemadeAdqusicionEMG/EMG_desarrollo/base_de_datos_electrodos/2026-09-22/O_Prueba8_Sujeto1/plot_calibrado_2026-09-22_O_Prueba8_Sujeto1.png}
\caption{Señales calibradas de la vocal O - O\_Prueba8\_Sujeto1: filtro Notch en 50 Hz, pasabanda y envolvente RMS.}
\end{figure}

\subsubsection{Medición: O\_Prueba9\_Sujeto1}
\begin{figure}[H]
\centering
\includegraphics[width=0.88\textwidth]{/home/santiago/repositorios/Nandu_SistemadeAdqusicionEMG/EMG_desarrollo/base_de_datos_electrodos/2026-09-22/O_Prueba9_Sujeto1/plot_paper_combined.png}
\caption{Registro de 3 músculos de la vocal O - O\_Prueba9\_Sujeto1: línea temporal de ventanas y concatenación simétrica con ruido restado.}
\end{figure}

\begin{figure}[H]
\centering
\includegraphics[width=0.88\textwidth]{/home/santiago/repositorios/Nandu_SistemadeAdqusicionEMG/EMG_desarrollo/base_de_datos_electrodos/2026-09-22/O_Prueba9_Sujeto1/plot_calibrado_2026-09-22_O_Prueba9_Sujeto1.png}
\caption{Señales calibradas de la vocal O - O\_Prueba9\_Sujeto1: filtro Notch en 50 Hz, pasabanda y envolvente RMS.}
\end{figure}

\subsection{Vocal U}
Se presentan a continuación las tomas registradas para la vocal \textbf{U}, ordenadas cronológicamente por número de prueba:

\subsubsection{Medición: U\_Prueba6\_Sujeto1}
\begin{figure}[H]
\centering
\includegraphics[width=0.88\textwidth]{/home/santiago/repositorios/Nandu_SistemadeAdqusicionEMG/EMG_desarrollo/base_de_datos_electrodos/2026-09-22/U_Prueba6_Sujeto1/plot_paper_combined.png}
\caption{Registro de 3 músculos de la vocal U - U\_Prueba6\_Sujeto1: línea temporal de ventanas y concatenación simétrica con ruido restado.}
\end{figure}

\begin{figure}[H]
\centering
\includegraphics[width=0.88\textwidth]{/home/santiago/repositorios/Nandu_SistemadeAdqusicionEMG/EMG_desarrollo/base_de_datos_electrodos/2026-09-22/U_Prueba6_Sujeto1/plot_calibrado_2026-09-22_U_Prueba6_Sujeto1.png}
\caption{Señales calibradas de la vocal U - U\_Prueba6\_Sujeto1: filtro Notch en 50 Hz, pasabanda y envolvente RMS.}
\end{figure}

\subsubsection{Medición: U\_Prueba7\_Sujeto1}
\begin{figure}[H]
\centering
\includegraphics[width=0.88\textwidth]{/home/santiago/repositorios/Nandu_SistemadeAdqusicionEMG/EMG_desarrollo/base_de_datos_electrodos/2026-09-22/U_Prueba7_Sujeto1/plot_paper_combined.png}
\caption{Registro de 3 músculos de la vocal U - U\_Prueba7\_Sujeto1: línea temporal de ventanas y concatenación simétrica con ruido restado.}
\end{figure}

\begin{figure}[H]
\centering
\includegraphics[width=0.88\textwidth]{/home/santiago/repositorios/Nandu_SistemadeAdqusicionEMG/EMG_desarrollo/base_de_datos_electrodos/2026-09-22/U_Prueba7_Sujeto1/plot_calibrado_2026-09-22_U_Prueba7_Sujeto1.png}
\caption{Señales calibradas de la vocal U - U\_Prueba7\_Sujeto1: filtro Notch en 50 Hz, pasabanda y envolvente RMS.}
\end{figure}

\subsubsection{Medición: U\_Prueba8\_Sujeto1}
\begin{figure}[H]
\centering
\includegraphics[width=0.88\textwidth]{/home/santiago/repositorios/Nandu_SistemadeAdqusicionEMG/EMG_desarrollo/base_de_datos_electrodos/2026-09-22/U_Prueba8_Sujeto1/plot_paper_combined.png}
\caption{Registro de 3 músculos de la vocal U - U\_Prueba8\_Sujeto1: línea temporal de ventanas y concatenación simétrica con ruido restado.}
\end{figure}

\begin{figure}[H]
\centering
\includegraphics[width=0.88\textwidth]{/home/santiago/repositorios/Nandu_SistemadeAdqusicionEMG/EMG_desarrollo/base_de_datos_electrodos/2026-09-22/U_Prueba8_Sujeto1/plot_calibrado_2026-09-22_U_Prueba8_Sujeto1.png}
\caption{Señales calibradas de la vocal U - U\_Prueba8\_Sujeto1: filtro Notch en 50 Hz, pasabanda y envolvente RMS.}
\end{figure}

\subsubsection{Medición: U\_Prueba9\_Sujeto1}
\begin{figure}[H]
\centering
\includegraphics[width=0.88\textwidth]{/home/santiago/repositorios/Nandu_SistemadeAdqusicionEMG/EMG_desarrollo/base_de_datos_electrodos/2026-09-22/U_Prueba9_Sujeto1/plot_paper_combined.png}
\caption{Registro de 3 músculos de la vocal U - U\_Prueba9\_Sujeto1: línea temporal de ventanas y concatenación simétrica con ruido restado.}
\end{figure}

\begin{figure}[H]
\centering
\includegraphics[width=0.88\textwidth]{/home/santiago/repositorios/Nandu_SistemadeAdqusicionEMG/EMG_desarrollo/base_de_datos_electrodos/2026-09-22/U_Prueba9_Sujeto1/plot_calibrado_2026-09-22_U_Prueba9_Sujeto1.png}
\caption{Señales calibradas de la vocal U - U\_Prueba9\_Sujeto1: filtro Notch en 50 Hz, pasabanda y envolvente RMS.}
\end{figure}

\newpage
\section{Análisis Multimodal de Señales y Espectrogramas}
A continuación se presentan los registros multimodales de 4 paneles para las mediciones de la sesión, agrupados vocal por vocal. Cada figura integra el espectrograma acústico STFT con pre-énfasis, el oscilograma del micrófono rectificado con su envolvente acústica, la activación EMG normalizada por el Supremo Tricanal del pulso y el espectrograma muscular RGB:

\subsection{Vocal O: Espectrogramas y Registro Multimodal}

\begin{figure}[H]
\centering
\includegraphics[width=0.88\textwidth]{/home/santiago/repositorios/Nandu_SistemadeAdqusicionEMG/EMG_desarrollo/base_de_datos_electrodos/2026-09-22/O_Prueba6_Sujeto1/plot_espectrograma_multimodal.png}
\caption{Análisis multimodal de 4 paneles para la vocal O - O\_Prueba6\_Sujeto1: espectrograma acústico con pre-énfasis, señal de micrófono rectificada, activación muscular sEMG por Supremo Tricanal y espectrograma RGB.}
\end{figure}

\begin{figure}[H]
\centering
\includegraphics[width=0.88\textwidth]{/home/santiago/repositorios/Nandu_SistemadeAdqusicionEMG/EMG_desarrollo/base_de_datos_electrodos/2026-09-22/O_Prueba7_Sujeto1/plot_espectrograma_multimodal.png}
\caption{Análisis multimodal de 4 paneles para la vocal O - O\_Prueba7\_Sujeto1: espectrograma acústico con pre-énfasis, señal de micrófono rectificada, activación muscular sEMG por Supremo Tricanal y espectrograma RGB.}
\end{figure}

\begin{figure}[H]
\centering
\includegraphics[width=0.88\textwidth]{/home/santiago/repositorios/Nandu_SistemadeAdqusicionEMG/EMG_desarrollo/base_de_datos_electrodos/2026-09-22/O_Prueba8_Sujeto1/plot_espectrograma_multimodal.png}
\caption{Análisis multimodal de 4 paneles para la vocal O - O\_Prueba8\_Sujeto1: espectrograma acústico con pre-énfasis, señal de micrófono rectificada, activación muscular sEMG por Supremo Tricanal y espectrograma RGB.}
\end{figure}

\begin{figure}[H]
\centering
\includegraphics[width=0.88\textwidth]{/home/santiago/repositorios/Nandu_SistemadeAdqusicionEMG/EMG_desarrollo/base_de_datos_electrodos/2026-09-22/O_Prueba9_Sujeto1/plot_espectrograma_multimodal.png}
\caption{Análisis multimodal de 4 paneles para la vocal O - O\_Prueba9\_Sujeto1: espectrograma acústico con pre-énfasis, señal de micrófono rectificada, activación muscular sEMG por Supremo Tricanal y espectrograma RGB.}
\end{figure}

\subsection{Vocal U: Espectrogramas y Registro Multimodal}

\begin{figure}[H]
\centering
\includegraphics[width=0.88\textwidth]{/home/santiago/repositorios/Nandu_SistemadeAdqusicionEMG/EMG_desarrollo/base_de_datos_electrodos/2026-09-22/U_Prueba6_Sujeto1/plot_espectrograma_multimodal.png}
\caption{Análisis multimodal de 4 paneles para la vocal U - U\_Prueba6\_Sujeto1: espectrograma acústico con pre-énfasis, señal de micrófono rectificada, activación muscular sEMG por Supremo Tricanal y espectrograma RGB.}
\end{figure}

\begin{figure}[H]
\centering
\includegraphics[width=0.88\textwidth]{/home/santiago/repositorios/Nandu_SistemadeAdqusicionEMG/EMG_desarrollo/base_de_datos_electrodos/2026-09-22/U_Prueba7_Sujeto1/plot_espectrograma_multimodal.png}
\caption{Análisis multimodal de 4 paneles para la vocal U - U\_Prueba7\_Sujeto1: espectrograma acústico con pre-énfasis, señal de micrófono rectificada, activación muscular sEMG por Supremo Tricanal y espectrograma RGB.}
\end{figure}

\begin{figure}[H]
\centering
\includegraphics[width=0.88\textwidth]{/home/santiago/repositorios/Nandu_SistemadeAdqusicionEMG/EMG_desarrollo/base_de_datos_electrodos/2026-09-22/U_Prueba8_Sujeto1/plot_espectrograma_multimodal.png}
\caption{Análisis multimodal de 4 paneles para la vocal U - U\_Prueba8\_Sujeto1: espectrograma acústico con pre-énfasis, señal de micrófono rectificada, activación muscular sEMG por Supremo Tricanal y espectrograma RGB.}
\end{figure}

\begin{figure}[H]
\centering
\includegraphics[width=0.88\textwidth]{/home/santiago/repositorios/Nandu_SistemadeAdqusicionEMG/EMG_desarrollo/base_de_datos_electrodos/2026-09-22/U_Prueba9_Sujeto1/plot_espectrograma_multimodal.png}
\caption{Análisis multimodal de 4 paneles para la vocal U - U\_Prueba9\_Sujeto1: espectrograma acústico con pre-énfasis, señal de micrófono rectificada, activación muscular sEMG por Supremo Tricanal y espectrograma RGB.}
\end{figure}

\end{document}